\documentclass[11pt,a4paper,fontset=fandol,scheme=plain]{ctexart}

\usepackage[margin=25mm]{geometry}
\usepackage{microtype}
\usepackage{enumitem}
\usepackage{xcolor}
\usepackage{hyperref}

\definecolor{accent}{HTML}{285FAE}

\hypersetup{
  unicode=true,
  pdftitle={开篇},
  pdfauthor={Xayah Hina},
  pdfsubject={关于 i.xayah.me 写作空间的介绍},
  colorlinks=true,
  linkcolor=accent,
  urlcolor=accent
}

\setlength{\parindent}{2em}
\setlength{\parskip}{0pt}
\setlist{leftmargin=1.5em,itemsep=0.4em,topsep=0.5em}
\setcounter{secnumdepth}{3}

\title{开篇}
\author{Xayah Hina}
\date{2026-07-15}

\begin{document}
\maketitle

\begin{abstract}
一个关于图形学、音乐、技术工作与普通日常的写作空间，从这里开始。
\end{abstract}

\section{为什么有这个空间}

有些想法出现在技术工作的间隙里，还没有找到落脚的地方就消失了。另一些想法来自一段音乐、
一个视觉细节，或者某个普通得似乎不值得写、却又重要得不应该被忘记的瞬间。
这个网站想同时容纳这两类东西。

我的学术主页会尽量保持简洁，只放研究兴趣、论文、项目，以及那些可以被明确命名的工作。
这里则可以用另一种速度前进：记录尚未完成的问题，写下代码结束之后才出现的解释，
也保存那些可能要过很久才显出意义的观察。

\section{我想在这里写什么}

这些主题会不断变化，但大概总会回到几条熟悉的线索：

\begin{itemize}
  \item \textbf{图形学与研究。} 关于模拟、渲染、逆向设计、工具、失败的方法，
    以及那些很难放进项目页面里的推理过程。
  \item \textbf{音乐与聆听。} 记录我正在创作、听见，以及从声音中学到的东西。
  \item \textbf{日常生活。} 短小的随笔、旅途记录、用文字留下的照片，
    以及不必被包装成宏大结论也值得保存的时刻。
\end{itemize}

\section{关于形式}

每篇文章都会用 \LaTeX{} 写成，并以 PDF 发布。这个选择有意比在信息流里发一条内容更慢一些。
排版迫使我认真决定一个想法从哪里开始、在哪里结束；PDF 也让每篇文章拥有稳定的形态，
即使离开网站仍然可以被阅读和保存。

网页本身会保持为一个索引，而不是变成另一套复杂的发布系统。
它告诉你这里有什么，真正的文字则留在文档中。

\section{一个开始}

这篇开场有意保持简单。它只是标记一个起点，也为之后的文字留出足够的空白，
让这个网站的性格慢慢自然形成。如果这里的某些内容与你产生了共鸣，
也可以在 \href{https://graphics.xayah.me/}{graphics.xayah.me} 找到我的学术工作。

\end{document}
